% GENERATED FILE. Edit canonical lessons, then run npm run generate:books.
% canonical-source-hash: fnv1a64:2b58058dcceee1c2
% canonical-lessons: SA-C342-adharah, SA-C342-yuktih, SA-C342-arthat, SA-C342-phalatah, SA-C342-upapadayati

\chapter{On What Grounds}
\label{ch:sa-on-what-grounds}
\hlchaptermodality{342}

\begin{chapteropening}
\textbf{By the end of this chapter:} \emph{I can say what a view rests on and what follows from it.}

Complete the last lesson of chapter 342: I can say what a view rests on and what follows from it.
\end{chapteropening}

\section[\texorpdfstring{\sk{आधारः} (ādhāraḥ)}{आधारः (ādhāraḥ)}]{\sk{आधारः} (ādhāraḥ) — a basis; a support}

\label{lesson:SA-C342-adharah}



\begin{quote}
Before the new one: say the Sanskrit for urges on, then the Sanskrit for has something done.
\end{quote}

\subsection*{You'll want to know: \sk{आधारः}}
\textbf{\sk{आधारः}} — \emph{ādhāraḥ} — \textquotedblleft{}a basis; a support\textquotedblright{}.

\begin{quote}
\sk{तव} \sk{मतस्य} \sk{आधारः} \sk{कः}?
\end{quote}

\begin{quote}
\emph{tava matasya ādhāraḥ kaḥ?}
\end{quote}

\begin{quote}
What is the basis of your opinion?
\end{quote}

A masculine noun.

\begin{cousinweb}
\textbf{\sk{आधारः}} (\emph{ādhāraḥ}) is built on \emph{ā} and the root √dhṛ, \textquotedblleft{}to hold, to support\textquotedblright{}; MW gives \emph{ā-dhāra} the senses \textquotedblleft{}support, prop, substratum\textquotedblright{} (MW).
\end{cousinweb}

\begin{grammarlens}[title={What you've built}]
One more word for giving reasons and opinions.
\end{grammarlens}

\subsection*{Your turn}
\begin{itemize}
\raggedright
  \item \emph{Say it:} \emph{ādhāraḥ}
  \item \emph{Say it:} \emph{ādhāraḥ}, once more
  \item \emph{Say it:} \emph{tava matasya ādhāraḥ kaḥ?} — What is the basis of your opinion?
  \item \emph{Recall:} say the Sanskrit for urges on, then the Sanskrit for has something done, then say \emph{ādhāraḥ} again
\end{itemize}

\begin{tcolorbox}[breakable,colback=teal!4,colframe=teal!35!black,title={Before you move on}]
What does \sk{आधारः} mean? (\textquotedblleft{}A basis; a support\textquotedblright{}.) Say it once more.
\end{tcolorbox}
\section[\texorpdfstring{\sk{युक्तिः} (yuktiḥ)}{युक्तिः (yuktiḥ)}]{\sk{युक्तिः} (yuktiḥ) — a line of reasoning; an argument}

\label{lesson:SA-C342-yuktih}



\begin{quote}
Before the new one: say the Sanskrit for has something done, then the Sanskrit for a basis.
\end{quote}

\subsection*{You'll want to know: \sk{युक्तिः}}
\textbf{\sk{युक्तिः}} — \emph{yuktiḥ} — \textquotedblleft{}a line of reasoning; an argument\textquotedblright{}.

\begin{quote}
\sk{अरुणस्य} \sk{युक्तिः} \sk{दृढा} \sk{अस्ति।}
\end{quote}

\begin{quote}
\emph{aruṇasya yuktiḥ dṛḍhā asti.}
\end{quote}

\begin{quote}
Arun's argument is strong.
\end{quote}

A feminine noun.

\begin{cousinweb}
\textbf{\sk{युक्तिः}} (\emph{yuktiḥ}) comes from the root √yuj, \textquotedblleft{}to join\textquotedblright{}; MW gives \emph{yukti} the sense \textquotedblleft{}union\textquotedblright{} and then \textquotedblleft{}reasoning, argument, proof\textquotedblright{} (MW).
\end{cousinweb}

\begin{grammarlens}[title={What you've built}]
One more word for giving reasons and opinions.
\end{grammarlens}

\subsection*{Your turn}
\begin{itemize}
\raggedright
  \item \emph{Say it:} \emph{yuktiḥ}
  \item \emph{Say it:} \emph{yuktiḥ}, once more
  \item \emph{Say it:} \emph{aruṇasya yuktiḥ dṛḍhā asti.} — Arun's argument is strong.
  \item \emph{Recall:} say the Sanskrit for has something done, then the Sanskrit for a basis, then say \emph{yuktiḥ} again
\end{itemize}

\begin{tcolorbox}[breakable,colback=teal!4,colframe=teal!35!black,title={Before you move on}]
What does \sk{युक्तिः} mean? (\textquotedblleft{}A line of reasoning; an argument\textquotedblright{}.) Say it once more.
\end{tcolorbox}
\section[\texorpdfstring{\sk{अर्थात्} (arthāt)}{अर्थात् (arthāt)}]{\sk{अर्थात्} (arthāt) — that is to say}

\label{lesson:SA-C342-arthat}



\begin{quote}
Before the new one: say the Sanskrit for a basis, then the Sanskrit for a line of reasoning.
\end{quote}

\subsection*{You'll want to know: \sk{अर्थात्}}
\textbf{\sk{अर्थात्}} — \emph{arthāt} — \textquotedblleft{}that is to say\textquotedblright{}.

\begin{quote}
\sk{सः} \sk{कृपणः} \sk{अस्ति}, \sk{अर्थात्} \sk{धनं} \sk{न} \sk{ददाति।}
\end{quote}

\begin{quote}
\emph{saḥ kṛpaṇaḥ asti, arthāt dhanaṃ na dadāti.}
\end{quote}

\begin{quote}
He is stingy, that is to say he gives no money.
\end{quote}

\begin{cousinweb}
\textbf{\sk{अर्थात्}} (\emph{arthāt}) is built on \emph{artha}, \textquotedblleft{}aim, meaning\textquotedblright{}, in its ablative; MW gives \emph{arthāt} \textquotedblleft{}as a matter of fact\textquotedblright{} and \textquotedblleft{}according to the sense, that is to say\textquotedblright{} (MW).
\end{cousinweb}

\begin{grammarlens}[title={What you've built}]
One more word for giving reasons and opinions.
\end{grammarlens}

\subsection*{Your turn}
\begin{itemize}
\raggedright
  \item \emph{Say it:} \emph{arthāt}
  \item \emph{Say it:} \emph{arthāt}, once more
  \item \emph{Say it:} \emph{saḥ kṛpaṇaḥ asti, arthāt dhanaṃ na dadāti.} — He is stingy, that is to say he gives no money.
  \item \emph{Recall:} say the Sanskrit for a basis, then the Sanskrit for a line of reasoning, then say \emph{arthāt} again
\end{itemize}

\begin{tcolorbox}[breakable,colback=teal!4,colframe=teal!35!black,title={Before you move on}]
What does \sk{अर्थात्} mean? (\textquotedblleft{}That is to say\textquotedblright{}.) Say it once more.
\end{tcolorbox}
\section[\texorpdfstring{\sk{फलतः} (phalataḥ)}{फलतः (phalataḥ)}]{\sk{फलतः} (phalataḥ) — as a result}

\label{lesson:SA-C342-phalatah}



\begin{quote}
Before the new one: say the Sanskrit for a line of reasoning, then the Sanskrit for that is to say.
\end{quote}

\subsection*{You'll want to know: \sk{फलतः}}
\textbf{\sk{फलतः}} — \emph{phalataḥ} — \textquotedblleft{}as a result\textquotedblright{}.

\begin{quote}
\sk{सः} \sk{बहु} \sk{पठितवान्}, \sk{फलतः} \sk{परीक्षायां} \sk{प्रथमः} \sk{अभवत्।}
\end{quote}

\begin{quote}
\emph{saḥ bahu paṭhitavān, phalataḥ parīkṣāyāṃ prathamaḥ abhavat.}
\end{quote}

\begin{quote}
He studied hard; as a result he came first in the exam.
\end{quote}

\begin{cousinweb}
\textbf{\sk{फलतः}} (\emph{phalataḥ}) is built on \emph{phala}, \textquotedblleft{}fruit, result\textquotedblright{}, with the ending \emph{-tas}, \textquotedblleft{}from\textquotedblright{}; MW gives \emph{phala-tas} \textquotedblleft{}in relation to the reward or result\textquotedblright{} and \textquotedblleft{}consequently, accordingly\textquotedblright{} (MW).
\end{cousinweb}

\begin{grammarlens}[title={What you've built}]
One more word for giving reasons and opinions.
\end{grammarlens}

\subsection*{Your turn}
\begin{itemize}
\raggedright
  \item \emph{Say it:} \emph{phalataḥ}
  \item \emph{Say it:} \emph{phalataḥ}, once more
  \item \emph{Say it:} \emph{saḥ bahu paṭhitavān, phalataḥ parīkṣāyāṃ prathamaḥ abhavat.} — He studied hard; as a result he came first in the exam.
  \item \emph{Recall:} say the Sanskrit for a line of reasoning, then the Sanskrit for that is to say, then say \emph{phalataḥ} again
\end{itemize}

\begin{tcolorbox}[breakable,colback=teal!4,colframe=teal!35!black,title={Before you move on}]
What does \sk{फलतः} mean? (\textquotedblleft{}As a result\textquotedblright{}.) Say it once more.
\end{tcolorbox}
\section[\texorpdfstring{\sk{उपपादयति} (upapādayati)}{उपपादयति (upapādayati)}]{\sk{उपपादयति} (upapādayati) — demonstrates; argues for}

\label{lesson:SA-C342-upapadayati}



\begin{quote}
Before the new one: say the Sanskrit for that is to say, then the Sanskrit for as a result.
\end{quote}

\subsection*{You'll want to know: \sk{उपपादयति}}
\textbf{\sk{उपपादयति}} — \emph{upapādayati} — \textquotedblleft{}demonstrates; argues for\textquotedblright{}.

\begin{quote}
\sk{अधिवक्ता} \sk{प्रमाणेन} \sk{पक्षम्} \sk{उपपादयति।}
\end{quote}

\begin{quote}
\emph{adhivaktā pramāṇena pakṣam upapādayati.}
\end{quote}

\begin{quote}
The lawyer argues his side with evidence.
\end{quote}

Past (he …): \sk{सः} \sk{मतम्} \sk{उपपादितवान्}, \emph{saḥ matam upapāditavān}.

\begin{cousinweb}
\textbf{\sk{उपपादयति}} (\emph{upapādayati}) is built on \emph{upa}, \textquotedblleft{}towards\textquotedblright{}, and the root √pad, \textquotedblleft{}to go\textquotedblright{}, in its causative form; MW gives \emph{upa-√pad} the senses \textquotedblleft{}to accomplish, to produce\textquotedblright{} and \textquotedblleft{}to prove, to justify\textquotedblright{} (MW).
\end{cousinweb}

\begin{grammarlens}[title={What you've built}]
One more word for giving reasons and opinions.
\end{grammarlens}

\subsection*{Your turn}
\begin{itemize}
\raggedright
  \item \emph{Say it:} \emph{upapādayati}
  \item \emph{Say it:} \emph{upapādayati}, once more
  \item \emph{Say it:} \emph{adhivaktā pramāṇena pakṣam upapādayati.} — The lawyer argues his side with evidence.
  \item \emph{Recall:} say the Sanskrit for that is to say, then the Sanskrit for as a result, then say \emph{upapādayati} again
\end{itemize}

\begin{tcolorbox}[breakable,colback=teal!4,colframe=teal!35!black,title={Before you move on}]
What does \sk{उपपादयति} mean? (\textquotedblleft{}Demonstrates; argues for\textquotedblright{}.) Say it once more.
\end{tcolorbox}
